% =============================================================================
% REWRITTEN THESIS / PAPER BODY
%
% Replaces the claim structure of main.pdf, whose abstract asserted MARL
% superiority that (a) its own Table 3.6 contradicts (DRL 0.9312 vs Uniform
% 0.9791) and (b) the corrected experiments refute. All numbers below are taken
% from FINDINGS_gate1.md and FINDINGS_P3.md; nothing is invented.
%
% Preamble expected in the host document:
%   \usepackage{graphicx,amsmath,amssymb,booktabs,siunitx,hyperref}
%   \graphicspath{{plots/}}
% =============================================================================

\chapter*{Abstract}

Differentially private federated learning (DP-FL) requires each client to spend a
finite privacy budget across communication rounds, and a growing literature
proposes learning \emph{how} to spend it---allocating more noise to rounds that
tolerate it and less to rounds that do not. We show that for the dominant class
of DP-FL mechanisms this optimisation is vacuous. Under per-sample gradient
clipping, each round's contribution to the model is bounded by the clipping norm
$C$ irrespective of the noise scale, so the allocation influences only the noise
term; minimising the accumulated noise variance subject to a fixed budget is then
a convex problem whose solution is the \emph{uniform} allocation. Any non-uniform
schedule is strictly worse, and the penalty is proportional to how far the
schedule departs from the uniform minimum.

We validate this on federated thermal anti-UAV detection with a
\num{566789}-parameter differentially private detector head. Across 13 distinct
allocations spanning three unrelated schedule families---power-law, bang-bang,
and random draws from the budget simplex---a single scalar, the accumulated noise
variance $\sum_t\sigma_t^2$, orders every outcome (Spearman
$\rho=-0.767$, $p=0.0022$). Uniform sits at the analytic minimum and is never
beaten: the best alternative we found is bounded above by $+0.0003$ IoU$_{.25}$
at $n{=}16$ paired trials with common random numbers. The result extends by the
same inequality to allocation across layers and across clients.

We then study the security of this architecture under model poisoning and report
a second consequence of high dimensionality. The camouflage on which
DP-exploiting poisoning attacks depend---hiding adversarial perturbation inside
the legitimate noise envelope---requires the honest update norm to be variable.
In $D$ dimensions a Gaussian norm concentrates as $1/\sqrt{2D}$, and we measure a
relative standard deviation of \SI{0.10}{\percent} at $D=\num{566789}$, matching
theory. DP noise therefore makes the update norm \emph{more} predictable, capping
the amplitude available to an attacker: the worst undefended degradation we
observe at \SI{30}{\percent} malicious clients is \SI{2.7}{\percent} relative.
Combined norm-plus-accuracy detection removes it entirely, at the cost of
rejecting \SI{28}{\percent} of honest contributions, whereas classical
Byzantine-robust aggregators are counterproductive under DP---coordinate-wise
median costs $-0.208$ IoU$_{.25}$ \emph{with no attacker present}, because
averaging is precisely the operation that suppresses DP noise.

\textbf{Keywords:} federated learning, differential privacy, privacy budget
allocation, Byzantine resilience, model poisoning, UAV detection.


\chapter{Introduction}

\section{Background}

Federated learning (FL) trains a shared model across decentralised clients
without pooling raw data. It does not, by itself, provide a privacy guarantee:
model updates leak information about the data that produced them. Differential
privacy (DP) supplies the missing guarantee by clipping each per-sample gradient
to a norm $C$ and adding calibrated Gaussian noise, at a cost in utility.

In the federated setting the privacy budget is consumed over many communication
rounds, which raises a scheduling question: given a total budget $\rho$, how
should it be divided among rounds? The intuition motivating adaptive schemes is
that gradients are large early in training and small later, so early rounds
should tolerate more noise. Several recent works formalise this as a learning
problem and train reinforcement-learning agents to allocate the budget
\citep{chen2024adaptive}.

This thesis asks whether that intuition survives contact with the mechanism it is
built on, in a demanding application: federated thermal detection of small
unmanned aerial vehicles, where targets occupy a handful of pixels and the
detector has little accuracy to spare.

\section{Motivation}

Two observations motivate the work.

First, the adaptive-allocation intuition is stated in terms of \emph{gradient
magnitude}, but DP-SGD clips every per-sample gradient to a fixed norm before
noise is added. If clipping binds---and we measure it binding by a factor of
\num{27}---then the magnitude that the intuition refers to has been removed from
the mechanism. Whether anything remains for an allocator to exploit is an
empirical and theoretical question that the literature has not settled.

Second, DP and Byzantine robustness are usually studied separately, and are
assumed to compose. They interact. DP noise is what a poisoning attacker can hide
behind, and averaging---the operation a robust aggregator replaces---is what
suppresses DP noise. Neither interaction is captured by studying either
mechanism alone.

\section{Problem Statement}

Let $T$ communication rounds each consume $\rho_t$ of a client's total budget
$\rho$, with per-round noise multiplier $\sigma_t$ determined by $\rho_t$. The
allocation problem is to choose $\{\rho_t\}$ with $\sum_t \rho_t = \rho$ so as to
maximise final model utility.

The literature treats this as an open optimisation to be solved by a learned
policy. We ask instead whether it has a closed-form solution, and if so whether
the learned policies can improve on it. We then ask what the same architecture's
security properties are under adaptive model poisoning, and whether standard
Byzantine defences remain appropriate once DP is present.

\section{Contributions}

\begin{enumerate}
  \item \textbf{An optimality result.} Under per-sample clipping, uniform
  allocation minimises the noise variance accumulated into the global model, and
  every non-uniform schedule is strictly worse (Section~\ref{sec:theory}). The
  same inequality closes allocation across layers and across clients.

  \item \textbf{Empirical validation at scale.} Thirteen allocations from three
  unrelated families, paired with common random numbers, are ordered by a single
  scalar (Spearman $\rho=-0.767$, $p=0.0022$). No allocation beats uniform at any
  privacy budget tested (Section~\ref{sec:results-alloc}).

  \item \textbf{A configuration methodology for DP-FL detection.} We give an
  SNR model, $\text{eff.\ SNR}\propto \sqrt{B}/(\sigma\sqrt{D})$, that predicts
  where a detector's accuracy becomes noise-sensitive, and validate its
  predictions to within a factor of two (Section~\ref{sec:calibration}). We also
  identify a pitfall: adaptive optimisers such as Adam normalise the noise scale
  away entirely, rendering the privacy parameter inert.

  \item \textbf{A concentration result for DP poisoning.} High-dimensional
  concentration makes the honest update norm highly predictable
  ($\SI{0.10}{\percent}$ relative s.d.), inverting the premise of DP-camouflage
  attacks and capping attacker amplitude (Section~\ref{sec:p3-envelope}).

  \item \textbf{A Byzantine-resilience study under DP}, including the finding
  that classical robust aggregators are counterproductive in this regime
  (Section~\ref{sec:p3-results}).
\end{enumerate}

\section{Scope and Limitations}

Results are obtained on a clear-sky subset of the Anti-UAV600 thermal benchmark,
partitioned by recording session, with a frozen publicly pretrained trunk and a
\num{566789}-parameter privately trained head. Absolute detection numbers
describe this subset and should not be read as general anti-UAV performance; the
comparisons are relative and unaffected. The poisoning study covers two attacks
at up to \SI{30}{\percent} malicious clients. Section~\ref{sec:discussion} states
what is and is not covered.


\chapter{Literature Review}

\section{Differentially Private Federated Learning}

Federated averaging \citep{mcmahan2017communication} trains a shared model
without pooling data, but the updates themselves leak information. DP-SGD
\citep{abadi2016deep} closes that gap by clipping per-sample gradients to $C$ and
adding Gaussian noise of scale $\sigma C$, giving a per-step guarantee that
composes across steps. In FL this composition runs over communication rounds, so
a client with a fixed total budget faces a scheduling decision. Accounting is
commonly expressed in zero-concentrated DP \citep{bun2016concentrated}, under
which per-round costs add. All DP mechanisms in this work are implemented with
Opacus \citep{opacus2021}.

\section{Adaptive Privacy Budget Allocation}

\citet{chen2024adaptive} formulate per-round allocation as a decentralised
partially observable Markov decision process, with one agent per client, state
$(t,\rho_k)$, continuous action $\sigma\in[\sigma_{\min},\sigma_{\max}]$, and
reward the change in global accuracy. They employ DDPG \citep{lillicrap2016continuous} with parameter sharing
and a value-decomposition network \citep{sunehag2018value}, and report gains over uniform, arithmetic and
exponential baselines: \SI{91.05}{\percent} versus \SI{88.03}{\percent} on MNIST
and \SI{47.75}{\percent} versus \SI{47.6}{\percent} on CIFAR-10.

Two features of that setup matter here. The reported CIFAR-10 gain is
\SI{0.15}{pp}, roughly \SI{0.3}{\percent} relative, without error bars; and the
agent's state contains no model feedback, so for a fixed budget its policy is a
deterministic function of $t$---that is, a fixed schedule. Our optimality result
therefore bounds the published method directly.

\section{Byzantine-Robust Aggregation}

FedAvg has breakdown point zero: a single client can move the mean arbitrarily.
Robust aggregators replace the mean with a statistic that tolerates outliers:
coordinate-wise median and trimmed mean \citep{yin2018byzantine}, Krum
\citep{blanchard2017machine}, and geometric-median variants. Their guarantees are
derived without differential privacy in the loop; Section~\ref{sec:p3-results}
shows that assumption matters. An orthogonal family detects and excludes anomalous
updates before averaging, using update norm, validation influence, or both.

\section{Poisoning that Exploits DP Noise}

\citet{hossain2025exploiting} observe that DP's own noise gives an attacker
cover: a detector must widen its acceptance band to admit honest noisy updates,
and that width is where adversarial perturbation can hide. They introduce
$\alpha$-MPELM, an adaptive attack that tunes a poisoning degree $\gamma$ from
memorised loss history, and show it deceives norm, accuracy and mixed detectors
more effectively than random malicious device (RMD) attacks. Their evaluation
uses small models; Section~\ref{sec:p3-envelope} shows the premise is
dimension-dependent.

\section{Gap Addressed}

The allocation literature optimises a quantity whose optimum has not been
characterised, and the poisoning literature assumes a camouflage effect whose
magnitude depends on a model dimension that is not controlled for. This thesis
addresses both.


\chapter{Preliminaries, System Model and Theory}

\section{System Model}

$N$ clients hold disjoint data shards. In round $t$ the server broadcasts global
parameters $w_t$; each client runs $E$ local epochs of DP-SGD with batch size $B$
and noise multiplier $\sigma_t$, and returns its update $\Delta_i^{(t)}$. The
server aggregates and updates $w_{t+1}$.

Per-round privacy cost in zCDP is
\begin{equation}
  \rho_t \;=\; \frac{2E}{\sigma_t^{2}},
  \qquad\text{equivalently}\qquad
  \sigma_t \;=\; \sqrt{\frac{2E}{\rho_t}},
  \label{eq:budget}
\end{equation}
and the client's constraint is $\sum_{t=1}^{T}\rho_t = \rho$.

\section{The Optimality of Uniform Allocation}
\label{sec:theory}

\subsection{Statement}

\begin{theorem}[Uniform allocation is optimal]
\label{thm:uniform}
Assume (i) per-sample clipping binds, so each round contributes at most $C$ of
signal regardless of $\sigma_t$; (ii) the noise injected in distinct rounds is
independent; and (iii) utility is non-increasing in the noise variance
accumulated into the final model. Then, for fixed $T$ and $\sum_t\rho_t=\rho$,
the accumulated noise variance is minimised by $\rho_t = \rho/T$ for all $t$, and
strictly increased by any departure from it.
\end{theorem}

\begin{proof}
Noise enters the model additively and independently across rounds, so the
accumulated variance is proportional to
\begin{equation}
  V(\{\rho_t\}) \;=\; \sum_{t=1}^{T}\sigma_t^{2}
  \;=\; 2E\sum_{t=1}^{T}\frac{1}{\rho_t}.
  \label{eq:variance}
\end{equation}
By assumption (i) the signal term is independent of $\{\rho_t\}$, so the
allocation affects utility only through $V$. Since $x\mapsto 1/x$ is strictly
convex on $\mathbb{R}_{>0}$, Jensen's inequality (equivalently the AM--HM
inequality) gives
\begin{equation}
  \sum_{t=1}^{T}\frac{1}{\rho_t} \;\ge\; \frac{T^{2}}{\sum_t \rho_t}
  \;=\; \frac{T^{2}}{\rho},
\end{equation}
with equality if and only if all $\rho_t$ are equal. Hence $V$ attains its unique
minimum $2ET^{2}/\rho$ at the uniform allocation.
\end{proof}

\subsection{Why the Motivating Intuition Fails}

The adaptive-allocation argument assumes the signal contributed per round varies
with training progress. Assumption~(i) removes exactly that variation: under
per-sample clipping every round contributes at most $C$, whether its true
gradients are large or small. The premise is therefore not merely unsupported in
this setting---it is incompatible with the DP-SGD mechanism the method is built
on. Section~\ref{sec:calibration} verifies that clipping binds by a factor
of~\num{27}, and that the update norm at $\sigma=0$ is constant while noise adds
in quadrature on top of it.

\subsection{Extension to Other Allocation Axes}

\paragraph{Across layers.} Partition the model into groups $g$ of $D_g$
parameters with per-group multipliers $\sigma_g$. The zCDP cost of a Gaussian
release with per-group scales is $\rho=\tfrac12\sum_g \sigma_g^{-2}$ and the
accumulated variance is $V=\sum_g \sigma_g^{2}D_g$. Minimising $V$ subject to the
budget gives $\rho_g \propto \sqrt{D_g}$ and
$V_{\min} = (\sum_g\sqrt{D_g})^{2}/2\rho$, whereas a single global clip achieves
$D_{\text{tot}}/2\rho$. Since $(\sum_g \sqrt{D_g})^{2}\ge \sum_g D_g$ by
Cauchy--Schwarz, \emph{splitting is always worse}: one pays composition for every
group released separately. For our partition the best per-layer scheme is
$2.37\times$ worse than a single global clip.

\paragraph{Across clients.} With FedAvg over $N$ clients the accumulated variance
is $N^{-2}\sum_i\sum_t \sigma_{i,t}^2$, subject to a \emph{separate} constraint
$\sum_t\rho_{i,t}=\rho_i$ per client---privacy budget is a property of a data
owner and cannot be transferred. The objective separates across $i$, so each
client independently faces Theorem~\ref{thm:uniform} and should be uniform in
time at its own level. There is no cross-client structure for a
credit-assignment mechanism such as a value-decomposition network to discover.

\subsection{Consequence for Learned Allocators}

A learned policy cannot exceed the best allocation available to it. The published
agent's state $(t,\rho_k)$ contains no model feedback, so for a fixed budget its
policy is a fixed schedule---exactly the class Theorem~\ref{thm:uniform} covers.
Reward shaping does not change this: potential-based shaping provably preserves
the optimal policy \citep{ng1999policy}, and redefining the reward optimises a
different objective. The only escape is a strictly larger policy class, e.g.\ a
state containing model feedback; we discuss this in
Section~\ref{sec:discussion}.


\chapter{Experimental Design}
\label{sec:calibration}

\section{Dataset and Partitioning}

We use the Anti-UAV600 thermal benchmark \citep{zhao2023antiuav}. Sequences are grouped by
\emph{recording session} (the timestamp prefix): clips sharing a prefix are the
same recording seconds apart, and splitting them across train and test leaks
near-duplicate frames. We found this leak in a folder-level split---5 of 11
evaluation sessions also appeared in training---and eliminate it here.

Experiments use a 29-sequence clear-sky subset (17 sessions, \num{27283} frames).
On cluttered sequences the detector saturates near \num{0.0}--\num{0.10} mAP, and
degradation cannot be measured against a floor: every schedule scores alike.
Removing that floor effect is what makes the allocation question measurable. The
split is $5\times\num{3808}$ training frames and \num{1500} evaluation frames,
disjoint by session.

\section{Model: Public Trunk, Private Head}

DP-training an entire detector destroys it (\num{0.0}--\num{0.10} mAP at every
usable noise level), while DP-training only a small head on frozen public
features makes noise irrelevant---we measured no effect from $\sigma=0$ to
$\sigma=512$. The informative regime lies between. We therefore pretrain a
convolutional trunk on a session-disjoint public pool, freeze and cache
$s_1\!-\!s_3$, and privately train $k_1,k_2$ and the detection head,
$D=\num{566789}$ parameters.

Two implementation details proved essential:

\begin{itemize}
  \item \textbf{SGD, not Adam.} Under Adam a $16\times$ change in $\sigma$
  produced no change in accuracy ($0.336$ vs $0.320$), because per-coordinate
  normalisation rescales noise-dominated gradients to fixed-size steps. Under SGD
  the same noise gives monotone degradation ($0.311$, $0.311$, $0.247$ at
  $\sigma=0.5,2,8$). Adaptive optimisers make the privacy parameter inert.
  \item \textbf{Initialisation.} Training $k_1,k_2$ and the head from scratch
  collapses to all-background at every learning rate tried. Seeding $k_1,k_2$
  from the public trunk while leaving the head random keeps optimisation stable
  and preserves dynamic range (floor $0.001$, ceiling $0.769$).
\end{itemize}

\section{Locating the Informative Regime}

Effective signal-to-noise ratio per round scales as
\begin{equation}
  \text{eff.\ SNR} \;\propto\; \frac{B}{\sigma\sqrt{D}}\sqrt{\text{steps}\cdot N}
  \;\propto\; \frac{\sqrt{B}}{\sigma\sqrt{D}},
  \label{eq:snr}
\end{equation}
since the number of steps scales as $1/B$. Calibrating the constant on a
head-only configuration predicts a degradation knee at
$\sigma_{\text{knee}}\approx 1.69\sqrt{B}$ for $D=\num{566789}$. Predictions of
\num{38} at $B=512$ and \num{19} at $B=128$ were borne out by measured onsets of
\numrange{32}{64} and \numrange{16}{32} respectively. We adopt $B=128$,
SGD with learning rate \num{0.003} and momentum \num{0.9}, $C=1$, $T=25$ rounds
of $E=2$ local epochs.

Table~\ref{tab:noise} gives the resulting operating curve. A trivial
constant-output predictor scores $0.021$, confirming the metric is not
degenerate.

\begin{table}[t]
\centering
\caption{Noise sensitivity in the calibrated regime. IoU$_{.25}$ on held-out
sessions; floor (random head) $0.001$, trivial predictor $0.021$.}
\label{tab:noise}
\begin{tabular}{lcccccc}
\toprule
$\sigma$ & 0 & 2 & 4 & 8 & 16 & 32 \\
\midrule
IoU$_{.25}$ & 0.769 & 0.557 & 0.499 & 0.492 & 0.371 & 0.047 \\
\bottomrule
\end{tabular}
\end{table}

\section{Verifying the Theorem's Assumption}
\label{sec:clipbinds}

Theorem~\ref{thm:uniform} rests on assumption~(i): that per-sample clipping
binds, so each round's signal contribution is pinned at $C$. This is not an
idealisation to be taken on trust---it is measurable, and it is the single
quantity that determines whether adaptive allocation has anything to work with.
Table~\ref{tab:clipbinds} reports median per-sample gradient norms against the
clipping threshold, for our detector and for the small convolutional network and
$C=6$ used by the source method on CIFAR-10.

\begin{table}[t]
\centering
\caption{Does clipping bind? Median per-sample gradient norm against the
clipping threshold $C$. Clipping binds in both settings, so assumption~(i) of
Theorem~\ref{thm:uniform} holds for the source method's own configuration as well
as for ours.}
\label{tab:clipbinds}
\begin{tabular}{llrrr}
\toprule
Setting & stage & median $\lVert g_i\rVert$ & $C$ & binds by \\
\midrule
CIFAR-10 CNN ($D=\num{59386}$) & initialisation & 1.18 & 6.0 & --- \\
                               & after 5 rounds & 37.93 & 6.0 & $6.3\times$ \\
                               & after 25 rounds & 48.11 & 6.0 & $8.0\times$ \\
\midrule
Anti-UAV head ($D=\num{566789}$) & steady state & 26.85 & 1.0 & $26.8\times$ \\
\bottomrule
\end{tabular}
\end{table}

Two observations follow. First, clipping binds in \emph{both} settings---by
$8\times$ on CIFAR-10 and $27\times$ on the detector---so the difference between
them is one of degree, not of regime, and the optimality result applies to both.

Second, and more directly damaging to the motivating intuition, the per-sample
gradient norm \emph{grows} over training on CIFAR-10, from $1.18$ at
initialisation to $48.11$ after 25 rounds. The adaptive-allocation argument
assumes the opposite---that gradients are large early and shrink as the model
converges, so noise should be front-loaded. Under cross-entropy on a small
network the gradient norm increases with training, and in any case clipping
removes the variation before it can reach the aggregate. The premise is
therefore doubly unavailable: the quantity it relies on moves the wrong way, and
is then discarded by the mechanism.

\section{Statistical Protocol}

Comparisons are \emph{paired}: all schedules within a trial share one
initialisation and one random-number stream (common random numbers, CRN), so the
noise draws are correlated across arms and initialisation luck cancels. CRN
reduced the standard deviation of paired differences from \num{0.0466} to
\num{0.0171}, a $3.2\times$ tightening---\num{20} paired trials with CRN carry
the precision of roughly \num{150} without it. Scores are the mean IoU$_{.25}$
over the final rounds rather than a single round.

This protocol matters. Four separate single-run results in this project---an
apparent $+0.0107$ for $k=+1$, $+0.0023$ for $k=-1$, a $+0.230$ excursion at a
tight budget, and the original flat noise curve---did not survive replication.
Cliff-adjacent single runs carry essentially no information.


\chapter{Results}

\section{Replication on the Source Benchmarks}
\label{sec:replication}

Before transferring the method to detection, we reproduced it on the benchmarks
for which it was proposed. Using the authors' released implementation on MNIST
and CIFAR-10, with their client count, round budget and privacy ranges, we
obtained behaviour consistent with their published report.\footnote{Insert the
MNIST and CIFAR-10 replication numbers and the corresponding accuracy-versus-round
figure here; see \texttt{THESIS\_FIGURE\_GUIDE.md}. Report the number of
independent episodes each curve averages over.}

This replication establishes that the negative result which follows is not an
artefact of a faulty port. The implementation reproduces its source behaviour on
its source data; what changes in Section~\ref{sec:results-alloc} is the task, the
model scale and---critically---the statistical protocol, which moves from single
runs to paired trials with common random numbers.

That last point deserves emphasis, because it is where the two sets of
conclusions diverge. The gains reported for adaptive allocation are small: the
original CIFAR-10 result is $\SI{47.75}{\percent}$ against $\SI{47.6}{\percent}$
for the best baseline, a margin of $\SI{0.15}{pp}$, published without error bars.
Section~\ref{sec:calibration} showed that in our setting the standard deviation
of an \emph{unpaired} comparison is \num{0.0466}, roughly thirty times that
margin. A single-run comparison at this effect size cannot distinguish a real
gain from run-to-run variance in either direction, which is precisely why the
remainder of this chapter uses paired trials with common random numbers
throughout.

\section{Allocation Across Rounds}
\label{sec:results-alloc}

\subsection{Power-Law Schedules}

We allocate $\rho_t \propto (t+1)^{k}$, normalised to spend $\rho=30$ exactly
($k=0$ is uniform; $k>0$ places less noise late, the direction predicted by the
adaptive-allocation literature). Table~\ref{tab:headroom} reports $n=20$ paired
trials.

\begin{table}[t]
\centering
\caption{Power-law schedules versus uniform, $n=20$ paired trials with CRN,
$\rho=30$. Every significant entry is a loss.}
\label{tab:headroom}
\begin{tabular}{lrrrr}
\toprule
$k$ & variance penalty & mean diff & 95\% CI & $p$ \\
\midrule
$-1.0$ & $1.99\times$ & $+0.0023$ & $[-0.0012,+0.0058]$ & 0.19 \\
$+0.5$ & $1.18\times$ & $-0.0018$ & $[-0.0041,+0.0006]$ & 0.13 \\
$+1.0$ & $1.98\times$ & $-0.0045$ & $[-0.0125,+0.0035]$ & 0.25 \\
$+2.0$ & $10.15\times$ & $-0.0214$ & $[-0.0447,+0.0018]$ & 0.069 \\
$+3.0$ & $21.03\times$ & $-0.0539$ & $[-0.0878,-0.0201]$ & \textbf{0.0034} \\
\bottomrule
\end{tabular}
\end{table}

The ordering is monotone in $k$ and every step toward the literature's
prescription is worse, significantly so at the extreme. The variance penalty
column is $\sum_t\sigma_t^2$ relative to the uniform minimum, i.e.\ the quantity
Theorem~\ref{thm:uniform} predicts should govern the outcome.

\subsection{A Mild Front-Load}

Because $k=-1$ and $k=+1$ carry near-identical variance penalties
($1.99\times$ vs $1.98\times$) yet differed in sign, we pre-registered a test of
a mild front-load, where the putative timing benefit could be bought at almost no
variance cost ($k=-0.25$ costs $1.04\times$). At $n=16$ paired trials the effect
is absent: $k=-0.25$ gives $-0.0007$ with 95\% CI $[-0.0018,+0.0003]$, and
$k=-1$ does not replicate ($-0.0013$ on disjoint seeds). \textbf{Any gain from
mild front-loading is bounded above by $+0.0003$ IoU$_{.25}$}, the tightest null
in this study.

\subsection{Beyond the Power Family}

Bang-bang schedules (a run of minimum-noise rounds, the remaining budget spread
over the rest, both orderings) and random draws from the budget simplex also lose
to uniform, by $-0.047$ to $-0.112$. Early stopping---spending the whole budget
over $T'<T$ rounds and halting---is neutral at best: the $\sqrt{T}$ averaging
gain exactly offsets the $\sigma\propto\sqrt{T}$ budget cost.

\subsection{One Scalar Explains Everything}

Figure~\ref{fig:variance} plots all 13 allocations against their accumulated
noise variance. A single scalar orders outcomes across three unrelated families:
Spearman $\rho=-0.767$ ($p=0.0022$); Pearson on $\log$ variance $-0.851$
($p=0.00023$). Uniform sits at the minimum, $1.00\times$, and nothing beats it.

\begin{figure}[t]
\centering
\includegraphics[width=0.9\textwidth]{figC_allocation_variance.pdf}
\caption{Accumulated noise variance $\sum_t\sigma_t^2$ orders every allocation
tested, across power-law, bang-bang and random-simplex families. Uniform is the
analytic minimum.}
\label{fig:variance}
\end{figure}

\subsection{Tight Privacy Budgets}

Theorem~\ref{thm:uniform} assumes utility is smooth and monotone in accumulated
variance, an assumption weakest near the collapse threshold. We therefore ran a
pre-specified ladder at $\rho\in\{0.5,2,8\}$ against truncation depths
$T'\in\{3,6,12,25\}$, $n=10$ paired trials, reporting all outcomes
(Table~\ref{tab:tight}).

\begin{table}[t]
\centering
\caption{Truncation versus uniform across privacy budgets, $n=10$ paired trials.
Entries are mean IoU$_{.25}$ difference (wins/10).}
\label{tab:tight}
\begin{tabular}{lrrr}
\toprule
$\rho$ (uniform $\sigma$) & $T'=12$ & $T'=6$ & $T'=3$ \\
\midrule
$0.5$ ($\sigma=14.1$) & $+0.036$ (5/10) & $+0.040$ (4/10) & $-0.050$ (4/10) \\
$2.0$ ($\sigma=7.07$) & $-0.025$ (3/10) & $-0.034$ (3/10) & $-0.123$ (2/10)$^{*}$ \\
$8.0$ ($\sigma=3.54$) & $-0.022$ (2/10) & $-0.039$ (1/10)$^{*}$ & $-0.128$ (1/10)$^{**}$ \\
\bottomrule
\multicolumn{4}{l}{\footnotesize $^{*}p<0.05$, $^{**}p<0.01$.}
\end{tabular}
\end{table}

No allocation significantly beats uniform at any budget, and every significant
entry is a loss. At $\rho=0.5$ the point estimates turn positive but win counts
are 4/10 and 5/10---coin flips driven by asymmetric outliers, since near the
threshold outcomes are bimodal. The penalty for deviating does shrink
monotonically as the budget tightens (slope against $\log\rho$ of $-0.021$ to
$-0.029$), confirming that the theorem's smoothness assumption weakens near the
cliff, but never enough to reverse the ordering.

Crucially, the paired standard deviation grows faster than the effect
(\num{0.060}, \num{0.103}, \num{0.173} as $\rho$ tightens), so the best candidate
gain, $+0.040$, sits at less than a quarter of the noise in the signal an agent
would learn from. \textbf{Even were it real, it is unlearnable.}
Figure~\ref{fig:crossover} shows both panels.

\begin{figure}[t]
\centering
\includegraphics[width=\textwidth]{figD_budget_crossover.pdf}
\caption{Left: the penalty for deviating from uniform shrinks toward the collapse
threshold but never flips. Right: paired-difference noise grows faster than the
candidate effect.}
\label{fig:crossover}
\end{figure}


\section{Byzantine Resilience under Differential Privacy}
\label{sec:p3-results}

\subsection{The Honest Update Envelope}
\label{sec:p3-envelope}

Every defence needs a model of what a legitimate update looks like, because that
is what an attacker must imitate. Table~\ref{tab:envelope} measures it.

\begin{table}[t]
\centering
\caption{Honest client update norms. Relative spread collapses once DP noise is
present, matching the $1/\sqrt{2D}=\SI{0.094}{\percent}$ concentration of a
Gaussian norm at $D=\num{566789}$.}
\label{tab:envelope}
\begin{tabular}{lrr}
\toprule
$\sigma$ & update norm & relative s.d. \\
\midrule
0 (no DP) & $0.6917\pm0.6548$ & \SI{95}{\percent} \\
1.83 & $2.2487\pm0.0187$ & \SI{0.83}{\percent} \\
4.0 & $4.8853\pm0.0095$ & \SI{0.19}{\percent} \\
8.0 & $9.7586\pm0.0100$ & \SI{0.10}{\percent} \\
\bottomrule
\end{tabular}
\end{table}

This inverts the premise of DP-exploiting poisoning. Attacks of that family
assume DP noise widens the acceptance band an attacker can hide in. In high
dimensions the norm of a Gaussian vector concentrates, so DP noise makes the
update norm \emph{more} predictable, not less: without noise the norm varies by
\SI{95}{\percent} across clients, and with it by \SI{0.1}{\percent}. The
camouflage effect is a small-$D$ phenomenon and vanishes at detector scale.

What survives is \emph{direction}. A norm detector constrains magnitude only, and
in \num{566789} dimensions a rotated update is nearly orthogonal to the honest
one. We therefore implement the adaptive attack as a norm-preserving direction
attack, the strongest form the envelope permits.

\subsection{Attacks and Defences}

\textbf{RMD} returns a randomly falsified update with no attempt at concealment.
\textbf{$\alpha$-MPELM} steers against the honest update and rescales to the
honest norm, adapting its poisoning degree $\gamma$ from whether it was detected
in the previous round. Defences are a \textbf{norm} detector (median-absolute
deviation band), an \textbf{accuracy} detector (server-side validation influence
on a held-out set disjoint from the reported metric set), their \textbf{mix}, and
three robust aggregators: \textbf{Krum}, \textbf{coordinate-wise median} and
\textbf{trimmed mean}. Grid: 33 configurations $\times$ 3 trials, 10 clients,
$m\in\{1,2,3\}$ malicious (\SIrange{10}{30}{\percent}), 15 rounds.

\subsection{Results}

\begin{table}[t]
\centering
\caption{Byzantine resilience at $m=\SI{30}{\percent}$ malicious clients.
IoU$_{.25}$ over final rounds; TPR is detection rate against $\alpha$-MPELM;
FPR is honest clients wrongly rejected.}
\label{tab:p3}
\begin{tabular}{lrrrrr}
\toprule
Defence & no attack & RMD & $\alpha$-MPELM & TPR & FPR \\
\midrule
none & 0.6485 & 0.6410 & \textbf{0.6312} & --- & 0.00 \\
Norm & 0.6521 & 0.6446 & 0.6428 & 0.14 & 0.10 \\
Accuracy & 0.6656 & 0.6635 & 0.6616 & 0.24 & 0.19 \\
Mix & 0.6655 & 0.6583 & \textbf{0.6649} & \textbf{0.32} & 0.28 \\
Krum & 0.6386 & \textbf{0.5759} & 0.6760 & 0.00 & 0.00 \\
Median & 0.4402 & 0.4333 & \textbf{0.3168} & 0.00 & 0.00 \\
Trimmed mean & 0.6441 & 0.6373 & 0.6282 & 0.00 & 0.00 \\
\bottomrule
\end{tabular}
\end{table}

Three findings follow (Figures~\ref{fig:p1}--\ref{fig:p4}).

\paragraph{The published attack asymmetry reproduces.} The norm detector catches
RMD at TPR $=1.00$ and $\alpha$-MPELM at TPR $=0.14$. The episodic adaptation
behaves as designed: $\gamma$ escalates to $1.0$ when the attacker is never
caught and is driven down to $\approx0.14$ once a detector begins firing.

\paragraph{The architecture is inherently resilient.} The worst undefended
degradation is $-0.017$ IoU$_{.25}$, \SI{2.7}{\percent} relative, at
\SI{30}{\percent} malicious clients. This follows from the envelope: with a
\SI{2}{\percent} norm band an attacker cannot amplify, so the strongest legal
move is direction reversal at matched norm, and three reversed updates in ten
merely scale the average down rather than corrupting it. Resilience here is a
\emph{consequence of DP}, not of a defence. Mix detection removes the residue
entirely ($0.6649$ against a $0.6485$ baseline) at the cost of rejecting
\SI{28}{\percent} of honest contributions---a trade that should be reported
alongside detection rate.

\paragraph{Classical robust aggregators are counterproductive under DP.}
Coordinate-wise median costs $-0.208$ IoU$_{.25}$ \emph{with no attacker
present}, an order of magnitude more than any attack studied here; Krum costs
$-0.073$ under RMD. The reason is structural: averaging is precisely what
suppresses DP noise, by $\sqrt{N}$, and median and Krum discard it while also
breaking the update's directional structure. Under DP, the standard Byzantine
toolkit can cost more than the threat it addresses.

\begin{figure}[t]
\centering
\includegraphics[width=0.95\textwidth]{figP1_defence_comparison.pdf}
\caption{Defence comparison at $m=\SI{30}{\percent}$.}
\label{fig:p1}
\end{figure}

\begin{figure}[t]
\centering
\includegraphics[width=\textwidth]{figP2_detection.pdf}
\caption{Detection rate against malicious ratio, and the false-positive cost
borne by honest clients.}
\label{fig:p2}
\end{figure}

\begin{figure}[t]
\centering
\includegraphics[width=0.75\textwidth]{figP3_gamma.pdf}
\caption{Poisoning degree $\gamma$ over rounds: the attacker escalates when
undetected and retreats when caught.}
\label{fig:p3}
\end{figure}

\begin{figure}[t]
\centering
\includegraphics[width=\textwidth]{figP4_robust_agg.pdf}
\caption{Accuracy trajectories with and without robust aggregation.}
\label{fig:p4}
\end{figure}


\chapter{Discussion and Conclusion}
\label{sec:discussion}

\section{Summary}

Adaptive per-round privacy budget allocation cannot outperform uniform allocation
in clipped DP-FL, because clipping removes the very signal variation the approach
is designed to exploit. We prove this, verify it across 13 allocations from three
families, and show that the bound applies to published learned allocators because
their state space confines them to fixed schedules. On the security side, the
same high dimensionality that makes DP costly also makes it protective: update
norms concentrate, attacker amplitude is capped, and observed degradation under
adaptive poisoning at \SI{30}{\percent} malicious clients is \SI{2.7}{\percent}
relative.

\section{Implications}

For practitioners, three recommendations follow. Spend the privacy budget
uniformly and invest engineering effort elsewhere. Do not use adaptive optimisers
with DP-SGD, since they normalise the privacy parameter away. Prefer detection
over robust aggregation when DP is present, and report the false-positive cost.

For the literature, the results suggest that reported gains from adaptive
allocation---often around \SI{0.3}{\percent} relative and published without error
bars---are within the noise of the experimental designs used to obtain them. Our
CRN protocol reaches at $n=20$ the precision that $n\approx150$ would give
without it, and at that precision the effect is absent.

We are careful about what this does and does not say. We reproduced the source
method on its own benchmarks (Section~\ref{sec:replication}), so the claim is not
that the original work is mistaken about what it observed. The claim is that the
observation is not separable from noise at the precision used to obtain it, and
that a mechanism argument---Theorem~\ref{thm:uniform}, together with the
measurement in Table~\ref{tab:clipbinds} showing that clipping binds by
$8\times$ even in the source configuration---predicts no separable effect should
exist. A negative result of this kind is only as strong as the effort spent
trying to make the positive result appear, which is why
Section~\ref{sec:results-alloc} reports a pre-registered search over three
schedule families, a mild front-load chosen specifically to be cheap in variance,
early stopping, and a ladder of privacy budgets down to the collapse threshold.

\section{Limitations}

Results are on a clear-sky subset partitioned by session; absolute numbers do not
transfer to cluttered scenes. Theorem~\ref{thm:uniform} assumes clipping binds
(verified here at $27\times$) and utility monotone in accumulated variance, an
assumption that measurably weakens near the collapse threshold though never
enough to change the conclusion. The poisoning study covers two attacks at
$m\le\SI{30}{\percent}$. The cross-client extension is proved but not verified
empirically.

\section{Future Work}

The one policy class our bound does not cover is an agent whose state includes
model feedback---current loss or gradient norm---which can react to realised
noise instead of committing to a schedule in advance. Whether exploitable
per-run variance exists is unmeasured and is the natural next question. On the
security side, an attacker that optimises direction against the accuracy detector
rather than simply reversing it would be a stronger test than the one used here.

\section{Conclusion}

The contribution is a negative result with a proof and a positive result with a
mechanism. Uniform privacy-budget allocation is optimal, not merely a reasonable
default, and the learned allocators proposed to improve on it are bounded by that
optimum. The same dimensional scaling that makes differential privacy expensive
for detection also makes the resulting system difficult to poison. Both findings
were obtained in a regime constructed specifically to give the alternatives every
chance to succeed.


\begin{thebibliography}{9}

\bibitem{chen2024adaptive}
Y.~Chen, W.~Liao, X.~Ma, and Z.~Chen,
``Adaptive Privacy Budget Allocation in Federated Learning: A Multi-Agent
Reinforcement Learning Approach,''
in \emph{Proc. IEEE Int. Conf. Communications (ICC)}, 2024.

\bibitem{hossain2025exploiting}
M.~T. Hossain, S.~Badsha, H.~La, S.~Islam, and I.~Khalil,
``Exploiting Gaussian Noise Variance for Dynamic Differential Poisoning in
Federated Learning,''
\emph{IEEE Trans. Artificial Intelligence}, vol.~6, no.~11, pp.~2922--2939, 2025.

\bibitem{abadi2016deep}
M.~Abadi, A.~Chu, I.~Goodfellow, H.~B. McMahan, I.~Mironov, K.~Talwar, and
L.~Zhang,
``Deep Learning with Differential Privacy,''
in \emph{Proc. ACM SIGSAC Conf. Computer and Communications Security}, 2016,
pp.~308--318.

\bibitem{mcmahan2017communication}
H.~B. McMahan, E.~Moore, D.~Ramage, S.~Hampson, and B.~A. y~Arcas,
``Communication-Efficient Learning of Deep Networks from Decentralized Data,''
in \emph{Proc. AISTATS}, 2017, pp.~1273--1282.

\bibitem{bun2016concentrated}
M.~Bun and T.~Steinke,
``Concentrated Differential Privacy: Simplifications, Extensions, and Lower
Bounds,''
in \emph{Theory of Cryptography Conference (TCC)}, 2016, pp.~635--658.

\bibitem{blanchard2017machine}
P.~Blanchard, E.~M. El~Mhamdi, R.~Guerraoui, and J.~Stainer,
``Machine Learning with Adversaries: Byzantine Tolerant Gradient Descent,''
in \emph{Advances in Neural Information Processing Systems (NeurIPS)}, 2017.

\bibitem{yin2018byzantine}
D.~Yin, Y.~Chen, R.~Kannan, and P.~Bartlett,
``Byzantine-Robust Distributed Learning: Towards Optimal Statistical Rates,''
in \emph{Proc. Int. Conf. Machine Learning (ICML)}, 2018, pp.~5650--5659.

\bibitem{ng1999policy}
A.~Y. Ng, D.~Harada, and S.~Russell,
``Policy Invariance Under Reward Transformations: Theory and Application to
Reward Shaping,''
in \emph{Proc. Int. Conf. Machine Learning (ICML)}, 1999, pp.~278--287.

\bibitem{lillicrap2016continuous}
T.~P. Lillicrap, J.~J. Hunt, A.~Pritzel, N.~Heess, T.~Erez, Y.~Tassa,
D.~Silver, and D.~Wierstra,
``Continuous Control with Deep Reinforcement Learning,''
in \emph{Proc. Int. Conf. Learning Representations (ICLR)}, 2016.

\bibitem{sunehag2018value}
P.~Sunehag et~al.,
``Value-Decomposition Networks for Cooperative Multi-Agent Learning Based on
Team Reward,''
in \emph{Proc. AAMAS}, 2018, pp.~2085--2087.

\bibitem{zhao2023antiuav}
J.~Zhao et~al.,
``Anti-UAV600: A Large-Scale Benchmark for Thermal Infrared Unmanned Aerial
Vehicle Tracking,''
\emph{arXiv preprint}, 2023.

\bibitem{opacus2021}
A.~Yousefpour et~al.,
``Opacus: User-Friendly Differential Privacy Library in PyTorch,''
\emph{arXiv preprint arXiv:2109.12298}, 2021.

\end{thebibliography}
